\documentclass[12pt, titlepage]{article}

\usepackage[margin=1in]{geometry}
\usepackage{amsmath, amssymb}
\usepackage{booktabs}
\usepackage{tabularx}
\usepackage{enumitem}
\usepackage{xurl}
\usepackage{hyperref}
\hypersetup{colorlinks, linkcolor=red, urlcolor=blue}

\newcommand{\progname}{Natural Control}
\newcommand{\todo}[1]{\par\medskip\noindent\fbox{\parbox{0.95\linewidth}{\textbf{TODO (team):} #1}}\par\medskip}
\newcommand{\str}{\textit{string}}
\newcommand{\B}{\mathbb{B}}
\newcommand{\R}{\mathbb{R}}
\newcommand{\N}{\mathbb{N}}

% One access-program table: Name | In | Out | Exceptions
\newenvironment{programs}{\par\noindent\tabularx{\textwidth}{p{3.4cm}Xp{3.2cm}p{2.6cm}}\toprule\textbf{Name} & \textbf{In} & \textbf{Out} & \textbf{Exceptions}\\\midrule}{\bottomrule\endtabularx\par}

\begin{document}

\title{Module Interface Specification for \progname{}\\[4pt]\large Design Document Revision -1 (draft)}
\author{Team 7: Aaron Loh, Curtis Vos, Kevin Kim, Lucas DeBoer}
\date{November 11, 2026}
\maketitle

\section{Revision History}
\begin{tabularx}{\textwidth}{p{3cm}p{2cm}X}
\toprule
\textbf{Date} & \textbf{Version} & \textbf{Notes}\\
\midrule
2026-11-11 & 0.1 & First draft, from the simple POC. Syntax complete for every module; semantics brief, as Revision -1 asks.\\
\bottomrule
\end{tabularx}

\section{Symbols, Abbreviations and Acronyms}
See the SRS and the Module Guide (MG). Module numbers M1 to M12 are the MG's.

\section{Introduction}
This document gives the interface of every module in the MG: what each one exports, what state it keeps and what it assumes. Revision -1 concentrates on syntax; the semantics are short and will be completed in Revision 0.

\section{Notation}
The specification follows the notation of Hoffman and Strooper. Types used:

\begin{center}
\begin{tabular}{lll}
\toprule
\textbf{Type} & \textbf{Notation} & \textbf{Meaning}\\
\midrule
boolean & $\B$ & true or false\\
natural number & $\N$ & $0, 1, 2, \ldots$\\
real & $\R$ & a number of seconds or milliseconds\\
string & \str & a sequence of characters, such as a sentence\\
sequence & seq of $T$ & an ordered list of values of type $T$\\
optional & $T \cup \{\bot\}$ & a value of type $T$, or none ($\bot$)\\
\bottomrule
\end{tabular}
\end{center}

Types defined by modules: \textit{Command} (M4), \textit{State} and \textit{Task} (M5), \textit{Action}, \textit{Answer} and \textit{Decision} (M6), \textit{Email} and \textit{Event} (M9), \textit{StopTiming} (M11), \textit{Settings} (M12). \textit{DateTime} is a date and time of day. Out ``--'' means no output.

\section{Module Decomposition}
\begin{tabular}{p{0.3\textwidth} p{0.6\textwidth}}
\toprule
\textbf{Level 1} & \textbf{Level 2}\\
\midrule
Hardware-Hiding & M1 Voice\\
Behaviour-Hiding & M2 Accessible Client, M3 Supervisor, M4 Control Phrases, M5 Task State, M6 Action Gate, M7 Agent Tools\\
Software Decision & M8 AI Model, M9 Mail and Calendar, M10 Web Reader, M11 Event Log, M12 Settings\\
\bottomrule
\end{tabular}

% ------------------------------------------------------------------
\newpage
\section{MIS of Control Phrases (M4)}
\subsection{Module} Control Phrases
\subsection{Uses} None.
\subsection{Syntax}
\subsubsection{Exported Constants}
\textit{Command} $=$ \{STOP, CANCEL, CONTINUE, YES, NO, OTHER\}\\
PHRASES: \textit{Command} $\to$ set of \str{} (the fixed phrases for STOP, CANCEL, CONTINUE and NO)
\subsubsection{Exported Access Programs}
\begin{programs}
classify & text: \str & \textit{Command} & --\\
\end{programs}
\subsection{Semantics}
\subsubsection{State Variables} None (library module).
\subsubsection{Environment Variables} None.
\subsubsection{Assumptions} Text is English, as transcribed by the speech-to-text service or typed.
\subsubsection{Access Routine Semantics}
\noindent classify(text):
\begin{itemize}
\item output: $out :=$ STOP if the normalized text contains ``stop working'' or equals a STOP phrase; otherwise the command whose phrase set contains the normalized text; otherwise YES if the text is a short run of yes words that begins with one; otherwise OTHER.
\item exception: none.
\end{itemize}
\subsubsection{Local Functions}
normalize: \str{} $\to$ seq of \str{}: lower-case words with punctuation and leading or trailing filler words removed.

% ------------------------------------------------------------------
\newpage
\section{MIS of Task State (M5)}
\subsection{Module} Task State
\subsection{Uses} None.
\subsection{Syntax}
\subsubsection{Exported Constants}
\textit{State} $=$ \{Working, AwaitingApproval, Executing, Pausing, Paused, Cancelled, Done\}\\
RUNNING $=$ \{Working, AwaitingApproval, Executing\}\\
ALLOWED: \textit{State} $\to$ set of \textit{State} (the moves of the state machine, SRS S.9 simplified for the POC)
\subsubsection{Exported Types}
\textit{Task} (abstract data type)
\subsubsection{Exported Access Programs}
\begin{programs}
new Task & request: \str & \textit{Task} & --\\
move & next: \textit{State} & -- & IllegalMove\\
running & -- & $\B$ & --\\
state & -- & \textit{State} & --\\
doneWork & -- & seq of \str & --\\
\end{programs}
\subsection{Semantics}
\subsubsection{State Variables}
request: \str; \quad state: \textit{State}; \quad done: seq of \str{}
\subsubsection{Environment Variables} None.
\subsubsection{Assumptions} One task is active per session (MG UC4).
\subsubsection{Access Routine Semantics}
\noindent new Task(r): transition: request, state, done $:=$ r, Working, $\langle\rangle$; output: $out :=$ self.\\[4pt]
\noindent move(next): transition: state $:=$ next; exception: $exc :=$ (next $\notin$ ALLOWED(state) $\Rightarrow$ IllegalMove).\\[4pt]
\noindent running(): output: $out :=$ state $\in$ RUNNING.

% ------------------------------------------------------------------
\newpage
\section{MIS of Action Gate (M6)}
\subsection{Module} Action Gate
\subsection{Uses} Control Phrases (M4), Task State (M5).
\subsection{Syntax}
\subsubsection{Exported Constants}
\textit{Answer} $=$ \{YES, NO, CHANGE, WITHDRAWN\}
\subsubsection{Exported Types}
\textit{Action} $=$ tuple of (kind: \str, details: \str{} $\to$ \str, readback: \str)\\
\textit{Decision} $=$ tuple of (answer: \textit{Answer}, note: \str)
\subsubsection{Exported Access Programs}
\begin{programs}
new Gate & speak: \str{} $\to$ -- & Gate & --\\
ask & task: \textit{Task}, action: \textit{Action} & \textit{Decision} & --\\
answer & text: \str & $\B$ & --\\
withdraw & -- & -- & --\\
pending & -- & \textit{Action} $\cup \{\bot\}$ & --\\
\end{programs}
\subsection{Semantics}
\subsubsection{State Variables}
waiting: \textit{Action} $\cup \{\bot\}$
\subsubsection{Environment Variables} Speech output, through speak.
\subsubsection{Assumptions} At most one action waits at a time.
\subsubsection{Access Routine Semantics}
\noindent ask(task, action): if task is not running, $out := \langle$WITHDRAWN, ``''$\rangle$. Otherwise: task moves to AwaitingApproval; waiting $:=$ action; the readback is spoken; the program waits until answer or withdraw is called. Then waiting $:= \bot$; on YES the task moves to Executing; on NO or CHANGE it moves to Working; $out :=$ the decision.\\[4pt]
\noindent answer(text): if waiting $= \bot$, $out :=$ false. Otherwise the decision is YES, NO or CHANGE (with note $=$ text) as classify(text) (M4) is YES, NO or anything else; $out :=$ true.\\[4pt]
\noindent withdraw(): if an action is waiting, its decision is WITHDRAWN.

% ------------------------------------------------------------------
\newpage
\section{MIS of Agent Tools (M7)}
\subsection{Module} Agent Tools
\subsection{Uses} Task State (M5), Action Gate (M6), Mail and Calendar (M9), Web Reader (M10).
\subsection{Syntax}
\subsubsection{Exported Constants}
UNTRUSTED: \str{} (the note placed before text from emails and websites)\\
PAUSED: \str{} (the result returned when the task may not act)
\subsubsection{Exported Access Programs}
\begin{programs}
searchEmail & query: \str & \str & --\\
readWebsite & url: \str, question: \str & \str & --\\
sendEmail & to: \str, subject: \str, body: \str & \str & --\\
addEvent & title: \str, start: \textit{DateTime}, end: \textit{DateTime}, place: \str & \str & --\\
writing & -- & $\B$ & --\\
calls & -- & $\N$ & --\\
\end{programs}
\subsection{Semantics}
\subsubsection{State Variables}
inProgress: $\N$ (tool calls running); \quad sending: $\B$ (an approved send or add is running)
\subsubsection{Environment Variables} Speech output, for short progress lines.
\subsubsection{Assumptions} The current task is supplied by the Supervisor (M3).
\subsubsection{Access Routine Semantics}
\noindent searchEmail, readWebsite: if the task is not running, $out :=$ PAUSED; otherwise $out :=$ UNTRUSTED followed by the text found. No approval is needed (FR10).\\[4pt]
\noindent sendEmail, addEvent: if the task is not running, $out :=$ PAUSED. Otherwise the action and its readback go to ask (M6). On YES, the action runs to completion even if a stop arrives meanwhile, the result is looked up again (M9), and the work is recorded on the task; $out :=$ ``confirmed'' or ``outcome unknown''. On NO, CHANGE or WITHDRAWN, nothing is sent and $out$ says why.
\subsubsection{Local Functions}
spokenTime: \textit{DateTime} $\to$ \str{}: a time as it should sound, for example ``Saturday, November 21 at 9:00 AM''.

% ------------------------------------------------------------------
\newpage
\section{MIS of Supervisor (M3)}
\subsection{Module} Supervisor
\subsection{Uses} Control Phrases (M4), Task State (M5), Action Gate (M6), Agent Tools (M7), Event Log (M11).
\subsection{Syntax}
\subsubsection{Exported Constants}
SETTLE\_TIME: $\R = 5$ (seconds a pause waits for a send already in progress)\\
\textit{Key} $=$ \{stop, resume, cancel, approve, decline\}
\subsubsection{Exported Access Programs}
\begin{programs}
onPartial & text: \str & -- & --\\
onSentence & text: \str & $\B$ (true if handled here) & --\\
onKey & key: \textit{Key} & -- & --\\
onIdle & -- & -- & --\\
stop & source: \str & -- & --\\
resume & -- & -- & --\\
cancel & -- & -- & --\\
\end{programs}
\subsection{Semantics}
\subsubsection{State Variables}
task: \textit{Task} $\cup \{\bot\}$; \quad readbackHeard: $\B$; \quad lastFastStop: $\R$
\subsubsection{Environment Variables} Speech output; the AI's turn (interrupt and continue).
\subsubsection{Assumptions} Every finished sentence reaches onSentence before the AI sees it.
\subsubsection{Access Routine Semantics}
\noindent onSentence(text): STOP $\Rightarrow$ stop; CANCEL $\Rightarrow$ cancel; CONTINUE while paused $\Rightarrow$ resume; an action waiting $\Rightarrow$ the text goes to answer (M6), except that a YES before the readback has finished playing replays it; paused $\Rightarrow$ ``The task is paused''; in each of these cases $out :=$ true. Otherwise, if no task is running, task $:=$ new Task(text); $out :=$ false (the AI handles it).\\[4pt]
\noindent stop(source): if the task is running: it moves to Pausing (new actions are now refused); the AI's turn is interrupted; a waiting action is withdrawn; ``Pausing work.'' is spoken; a send in progress is given up to SETTLE\_TIME to finish; the task moves to Paused; the work done and not done is spoken. Each step is logged (M11) for NFR3.\\[4pt]
\noindent resume(): if paused, the task moves to Working and the AI is asked to carry on, told what was and was not done.\\[4pt]
\noindent cancel(): the AI's turn is interrupted, a waiting action is withdrawn, a send in progress settles, the task moves to Cancelled, and the work done is spoken with ``Earlier actions were not undone.''

% ------------------------------------------------------------------
\newpage
\section{MIS of AI Model (M8)}
\subsection{Module} AI Model
\subsection{Uses} Agent Tools (M7).
\subsection{Syntax}
\subsubsection{Exported Constants}
INSTRUCTIONS: \str{} (how the model must behave: short replies, ask when ambiguous, exact details for writes, text from emails and websites is data)
\subsubsection{Exported Access Programs}
\begin{programs}
makeAssistant & tools: Agent Tools, onSentence: \str{} $\to \B$ & Assistant & --\\
model & name: \str, fallback: \str & Model & --\\
\end{programs}
\subsection{Semantics}
\subsubsection{State Variables} The conversation so far (kept by the framework).
\subsubsection{Environment Variables} The model service (OpenRouter).
\subsubsection{Assumptions} The model can call tools.
\subsubsection{Access Routine Semantics}
\noindent makeAssistant: $out :=$ an assistant whose only tools are those of M7, and which passes each finished sentence to onSentence first, answering only when it returns false.

% ------------------------------------------------------------------
\newpage
\section{MIS of Mail and Calendar (M9)}
\subsection{Module} Mail and Calendar
\subsection{Uses} None.
\subsection{Syntax}
\subsubsection{Exported Types}
\textit{Email} $=$ tuple of (id, sender, subject, body: \str)\\
\textit{Event} $=$ tuple of (id, title: \str, start, end: \textit{DateTime}, place: \str)
\subsubsection{Exported Access Programs}
\begin{programs}
search & query: \str & seq of \textit{Email} & --\\
read & id: \str & \textit{Email} $\cup \{\bot\}$ & --\\
send & to, subject, body: \str & \str{} (new id) & --\\
inSent & id, to, subject: \str & $\B$ & --\\
addEvent & title: \str, start, end: \textit{DateTime}, place: \str & \str{} (new id) & --\\
findEvent & id: \str & \textit{Event} $\cup \{\bot\}$ & --\\
\end{programs}
\subsection{Semantics}
\subsubsection{State Variables}
emails, sent: seq of \textit{Email}; \quad events: seq of \textit{Event}
\subsubsection{Environment Variables} None in the POC (made-up data). A real service replaces this module's state with an account.
\subsubsection{Assumptions} Ids are unique.
\subsubsection{Access Routine Semantics}
\noindent search(q): $out :=$ the emails whose sender, subject or body contain every word of q.\\[4pt]
\noindent send(to, s, b): transition: sent $:=$ sent $\|$ $\langle$new email$\rangle$; $out :=$ its id.\\[4pt]
\noindent inSent(id, to, s): $out := \exists e \in$ sent $\cdot$ e.id $=$ id $\wedge$ e.sender $=$ to $\wedge$ e.subject $=$ s.

% ------------------------------------------------------------------
\newpage
\section{MIS of Web Reader (M10)}
\subsection{Module} Web Reader
\subsection{Uses} None.
\subsection{Syntax}
\subsubsection{Exported Constants}
ALLOWED\_ACTIONS $=$ \{open a page, follow a link, scroll, read\}
\subsubsection{Exported Access Programs}
\begin{programs}
readWebsite & url: \str, question: \str & \str & --\\
\end{programs}
\subsection{Semantics}
\subsubsection{State Variables} None.
\subsubsection{Environment Variables} A browser window; the network.
\subsubsection{Assumptions} The site is public; no sign-in is needed.
\subsubsection{Access Routine Semantics}
\noindent readWebsite(url, q): $out :=$ the answer to q found on the site, or a statement that it could not be read. Only ALLOWED\_ACTIONS are used, so nothing on the site is changed.

% ------------------------------------------------------------------
\newpage
\section{MIS of Event Log (M11)}
\subsection{Module} Event Log
\subsection{Uses} None.
\subsection{Syntax}
\subsubsection{Exported Types}
\textit{StopTiming} $=$ tuple of (blocked, acknowledged, paused: $\R$), each in milliseconds after the stop was heard, or $-1$ if that event never followed
\subsubsection{Exported Access Programs}
\begin{programs}
new EventLog & file: the log file & EventLog & --\\
write & event: \str, fields: set of (name: \str, value) & -- & --\\
stopTimings & log: the log file & seq of \textit{StopTiming} & --\\
\end{programs}
\subsection{Semantics}
\subsubsection{State Variables} logFile: the file this log appends to.
\subsubsection{Environment Variables} The log file; a monotonic clock.
\subsubsection{Assumptions} Callers never pass message text, speech or keys as fields (NFR5).
\subsubsection{Access Routine Semantics}
\noindent new EventLog(file): transition: logFile $:=$ file, creating its folder if it does not exist; output: $out :=$ self.\\[4pt]
\noindent write(e, f): transition: appends to logFile one record of e, f and the current monotonic time.\\[4pt]
\noindent stopTimings(log): $out :=$ one \textit{StopTiming} for each ``stop\_heard'' record s in log, in order. Each of its three values is the time from s to the first later record named ``blocked'', ``acknowledged'' or ``paused'' respectively, or $-1$ if there is no such record.

% ------------------------------------------------------------------
\newpage
\section{MIS of Settings (M12)}
\subsection{Module} Settings
\subsection{Uses} None.
\subsection{Syntax}
\subsubsection{Exported Types}
\textit{Settings} $=$ tuple of (model, fallbackModel, speechToText, textToSpeech: \str, realWebsites: $\B$, sendDelay, readDelay: $\R$)
\subsubsection{Exported Constants}
DEFAULT\_MODEL, DEFAULT\_FALLBACK\_MODEL: \str{} (the AI models used when none is set)\\
DEFAULT\_STT, DEFAULT\_TTS: \str{} (the speech services used when none is set)\\
By default realWebsites $=$ false and sendDelay $=$ readDelay $=$ 0.
\subsubsection{Exported Access Programs}
\begin{programs}
load & -- & \textit{Settings} & InvalidSetting\\
\end{programs}
\subsection{Semantics}
\subsubsection{State Variables} None.
\subsubsection{Environment Variables} The process environment (settings only; keys are read by the outside services, not by this module).
\subsubsection{Assumptions} None.
\subsubsection{Access Routine Semantics}
\noindent load(): output: $out :=$ the settings from the environment, each one that is not set taking its default; exception: $exc :=$ (sendDelay or readDelay is set but is not a number $\Rightarrow$ InvalidSetting).

% ------------------------------------------------------------------
\newpage
\section{MIS of Voice (M1)}
\subsection{Module} Voice
\subsection{Uses} Supervisor (M3), AI Model (M8), Event Log (M11), Settings (M12); it also creates Agent Tools (M7), Mail and Calendar (M9) and Web Reader (M10) for a session.
\subsection{Syntax}
\subsubsection{Exported Access Programs}
\begin{programs}
startSession & room: a LiveKit room & -- & --\\
say & text: \str, interruptible: $\B$ & speech handle & --\\
\end{programs}
\subsection{Semantics}
\subsubsection{State Variables} The voice session; the last task state shown to the client.
\subsubsection{Environment Variables} Microphone and speakers (through the client); the network; the speech-to-text and text-to-speech services.
\subsubsection{Assumptions} One user per room.
\subsubsection{Access Routine Semantics}
\noindent startSession(room): connects speech in and out; passes partial sentences to onPartial (M3), finished sentences through the assistant (M8) to onSentence, and key presses to onKey; publishes the task state to the client; calls onIdle once the AI, the tools and the speech have been quiet for two seconds.

% ------------------------------------------------------------------
\newpage
\section{MIS of Accessible Client (M2)}
\subsection{Module} Accessible Client
\subsection{Uses} Voice (M1), over the network.
\subsection{Syntax}
\subsubsection{Exported Access Programs}
The client's interface is the user's: each control below is a button with a keyboard shortcut.
\begin{programs}
startSession & -- & -- & ConnectionFailed\\
endSession & -- & -- & --\\
sendCommand & key: \textit{Key} (M3) & -- & --\\
sendText & text: \str & -- & --\\
\end{programs}
\subsection{Semantics}
\subsubsection{State Variables} connected: $\B$; \quad taskState: \textit{State} $\cup \{\bot\}$; \quad conversation: seq of \str
\subsubsection{Environment Variables} Screen, keyboard, microphone, speakers; VoiceOver.
\subsubsection{Assumptions} A current web browser.
\subsubsection{Access Routine Semantics}
\noindent sendCommand(key): sends key on the room's \texttt{control} topic. taskState follows the voice module's \texttt{task\_state} attribute and is announced when it changes.
\todo{Lucas: confirm the controls and shortcuts against src/web, and how state changes are announced.}

\newpage
\section{Appendix: Generative AI Use Disclosure}
\todo{Fill in as the rubric asks: the tools used and their versions, how and where each was used in this document, and how the content was checked. Each author adds their own part.}

\end{document}
